% TOP_PROBLEM: {"id": 30003895, "title": "All-Order Expansion of the Noncommutative $\\Phi^4$ Model", "category_id": 16, "category": {"id": 16, "name": "physics", "display_name": "Mathematical Physics", "description": "Problems at the intersection of mathematics and physics.", "slug": "physics", "order_index": 16, "created_at": "2026-07-31T15:26:25.671Z"}, "status": "open", "statusQualification": "The defining transform and boundary-value relations were recovered from PDF page 29. The report checked the conjectured coefficients through order ten and resummed the proposed series with Lagrange--B\\\"urmann inversion; that resummation alone does not establish that the actual nonlinear boundary-value solution has all those coefficients. The target remains coefficientwise formal validity, rather than an added convergence claim."}
% FIRST_POSTED: 2026-10-11
% ENTRY_KIND: statement_only
% SUBMISSION: {"kind": "statement_only", "category": "top_problems", "problem_id": "30003895", "proposal_type": "established", "source": "Non-commutative Geometry, Index Theory and Mathematical Physics (2018), Oberwolfach Reports, DOI 10.4171/owr/2018/32, pp. 1911--1981. Primary report accessed 2026-10-11; relevant PDF pages 28, 29, 30. https://doi.org/10.4171/owr/2018/32; Panzer--Wulkenhaar (catalogue-listed context; not independently consulted in this preparation). https://arxiv.org/abs/1807.02945; OWR report (catalogue-listed context; not independently consulted in this preparation). https://ivv5hpp.uni-muenster.de/u/raimar/publications/abstracts/owr18-raimar.pdf", "status_evidence": "Canonical source status open retained. Supplied literature triage dated 2026-08-21: The coefficient identity was checked through ten orders and then conjectured; the primary paper proves a Lambert-W resummation but still leaves the all-order combinatorial identity as a conjecture. Primary-report access checked 2026-10-11; this historical report and inherited triage do not certify current resolution of every component.", "status_checked_on": "2026-10-11", "duplicate_check": "Compared canonical ID, title and normalized target with the complete local UnsolvedMath registry; normalized exact-text matches: none; title matches: none. Related or overlapping canonical records: none identified in this batch. Primary source context and the existing OWR statement audit used for scope and duplicate review; an exhaustive mathematical-equivalence/literature search was not performed.", "attribution": "Original proposer not named in the supplied canonical record; attribution remains to the identified Oberwolfach contribution/report, not its editors as invented proposers. Explanatory statement preparation and source/semantic audit assisted by OpenAI GPT-6 Codex; no human mathematical authorship or independent proof verification is claimed."}
% SOURCE_STATEMENT: "For the function $I_\\lambda(a)$ arising from the boundary-value solution of the noncommutative $\\Phi^4$ model, does the formal expansion\n$$\nI_\\lambda(a)=\\sum_{n\\geq1}\\frac{\\lambda^n}{n!}\\frac{d^{n-1}}{da^{n-1}}\\bigl(-\\log(1+a)\\bigr)^n\n-\\lambda\\sum_{n\\geq1}\\frac{\\lambda^n}{n!}\\frac{d^{n-1}}{da^{n-1}}\\left[\\frac{\\bigl(-\\log(1+a)\\bigr)^n}{a}\\right]\n$$\nhold to every order in $\\lambda$?"
% Preparation and final self-audit: CONTRIBUTING.md prompts C and D; source/audit artifacts retained in the local physics_statements_2026_10_11 batch.
% Imported from: physics_statements.tex; UnsolvedMath ID 30003895
% !TeX program = lualatex
\documentclass[11pt]{article}
\usepackage[margin=25mm]{geometry}
\usepackage{fontspec}
\setmainfont{Latin Modern Roman}
\usepackage{amsmath,amssymb,mathtools}
\usepackage{unicode-math}
\setmathfont{Latin Modern Math}
\usepackage{xurl,hyperref}
\hypersetup{colorlinks=true,urlcolor=blue}
\setcounter{secnumdepth}{0}
\setlength{\parindent}{0pt}
\setlength{\parskip}{5pt}
\setlength{\emergencystretch}{3em}
\newcommand{\sref}[2]{\hyperlink{source:#1:#2}{[S#2]}}
\newcommand{\cataloguescope}[1]{\paragraph{Recorded target.}#1}
\newcommand{\R}{\mathbb R}
\newcommand{\C}{\mathbb C}
\newcommand{\N}{\mathbb N}
\newcommand{\Z}{\mathbb Z}
\newcommand{\Q}{\mathbb Q}
\newcommand{\F}{\mathbb F}
\newcommand{\eps}{\varepsilon}
\begin{document}
% UnsolvedMath ID 30003895; source catalogue code OWR-16407-007
\setcounter{section}{30003894}
\section{All-Order Expansion of the Noncommutative $\Phi^4$ Model}\label{30003895}
{\small UnsolvedMath ID 30003895 · Mathematical Physics}
\subsection{Definitions and mathematical statement}

\paragraph{Shared notation.}$\mathbb N=\{1,2,\ldots\}$; the standard number fields have their usual meanings. All additional conventions are stated locally.


\paragraph{Definitions and conventions.}
The coupling $\lambda$ is a formal expansion parameter and $a$ is the argument of the source\textquotesingle s boundary-value solution $I_\lambda$. Derivatives in the displayed series are ordinary derivatives with respect to $a$, and $(-\log(1+a))^n$ means the $n$th power before differentiation. Holding to every order means equality of each coefficient in the formal power series in $\lambda$, after the second sum\textquotesingle s extra factor $\lambda$ is taken into account; this does not by itself imply convergence for nonzero coupling. The primary definition uses the one-sided Hilbert transform $H_a[f]=\pi^{-1}\operatorname{p.v.}\int_0^{\Lambda^2}f(p)/(p-a)\,dp$ with a common cutoff, followed by its prescribed cutoff removal. Its angle functions satisfy $\lambda\pi\cot\tau_b(a)=1+a+b-\lambda\log a+I_\lambda(a)$ and $I_\lambda(a)=\pi^{-1}\int_0^\infty(e^{-H_p[\tau_a]}\sin\tau_a(p)-\lambda\pi/(1+p))\,dp$, in the source boundary-value convention for $a>0$.

\paragraph{Target retained from the canonical source catalogue.}
\paragraph{Catalogue wording retained for traceability.}
For the function $I_\lambda(a)$ arising from the boundary-value solution of the noncommutative $\Phi^4$ model, does the formal expansion
$$
I_\lambda(a)=\sum_{n\geq1}\frac{\lambda^n}{n!}\frac{d^{n-1}}{da^{n-1}}\bigl(-\log(1+a)\bigr)^n
-\lambda\sum_{n\geq1}\frac{\lambda^n}{n!}\frac{d^{n-1}}{da^{n-1}}\left[\frac{\bigl(-\log(1+a)\bigr)^n}{a}\right]
$$
hold to every order in $\lambda$?

\paragraph{Source scope and qualification.}
The defining transform and boundary-value relations were recovered from PDF page 29. The report checked the conjectured coefficients through order ten and resummed the proposed series with Lagrange--B\"urmann inversion; that resummation alone does not establish that the actual nonlinear boundary-value solution has all those coefficients. The target remains coefficientwise formal validity, rather than an added convergence claim.
\subsection{Short English statement}
Does the proposed derivative formula reproduce every perturbative coefficient of the noncommutative phi-four solution?
\subsection{Sources}
\begin{itemize}
\item[\hypertarget{source:30003895:1}{S1}] ULAM.AI and the UnsolvedMath Contributors, UnsolvedMath, record 30003895 (OWR-16407-007). Catalogue curation under CC BY 4.0; definitions and English exposition added locally.
\url{https://huggingface.co/datasets/ulamai/UnsolvedMath}.
\item[\hypertarget{source:30003895:2}{S2}] Non-commutative Geometry, Index Theory and Mathematical Physics (2018), Oberwolfach Reports, DOI 10.4171/owr/2018/32, pp. 1911--1981. Primary report accessed 2026-10-11; relevant PDF pages 28, 29, 30.
\url{https://doi.org/10.4171/owr/2018/32}.
\item[\hypertarget{source:30003895:3}{S3}] Panzer--Wulkenhaar (catalogue-listed context; not independently consulted in this preparation).
\url{https://arxiv.org/abs/1807.02945}.
\item[\hypertarget{source:30003895:4}{S4}] OWR report (catalogue-listed context; not independently consulted in this preparation).
\url{https://ivv5hpp.uni-muenster.de/u/raimar/publications/abstracts/owr18-raimar.pdf}.
\end{itemize}
\label{end:30003895}
\end{document}
